% Generated mechanism comparison; attributed source claims, not new proof claims.
\begin{longtable}{p{0.20\textwidth}p{0.32\textwidth}p{0.40\textwidth}}
\toprule Family & Designs & Distinguishing mechanism \\ \midrule\endhead
Original source / allocation & Primary PA194 (67/7); PA194 (68/8) & Same original command protocol and separate diagonal calculus. 67 readers minus seven folds, or 68 minus eight. \\
Constructor coefficient & M4R; M5; M6 & Coefficients four, five and six change numerical formula coding. Reassociated constructor decoder; equal total count. \\
Cached MP & PA194C; PA194CS & Successful MP stores the consequent in cache and clears visible top; Shift tests cache. CS reverses destructive comparison operands to remove restoration. \\
Visible no-push MP & PA194N; PA194N5; PA194N6 & Successful MP retains (b,w,b,0), omitting dummy-zero push; original sum guard. Coefficient variants retain 194 controls. \\
Cache-only guard & PA194H; PA194H5; PA194H6; PA194HS = filed PA194NS & Visible no-push MP with Shift test h=1; HS uses symmetric comparison. Four/five/six constructor variants and comparator sharing. \\
Symmetric sum guard & PA194NS (sum guard); PA194SUMS & Original sum condition t+h=1 with symmetric destructive comparison. Removes a source restoration loop without changing the guard. \\
Q distribution & PA194Q & Exchange the curried premises of the universal-distribution head. Recompiled 190-token emitter with cached controller. \\
Q with coefficient changes & PA194Q5; PA194Q6 & Q head plus coefficient five or six; original dummy-push MP. Different emitted numerical formulas at the same count. \\
Q with no-push MP & PA194QN = PA194NQ; PA194QN5; PA194QN6 & Q head, visible no-push MP, original sum guard. Four/five/six coefficient realizations. \\
Q with cache-only guard & PA194QH; PA194QH5; PA194QH6; PA194QHS & Q head, no-push MP, cache-only Shift. Coefficient variants; QHS uses symmetric comparison. \\
Q with symmetric sum guard & PA194QNS; PA194QS & Q head with the original sum guard; QNS removes MP push, QS retains it. Shared symmetric comparator. \\
Diagonal choice & PA194D; PA194DHS; Gamma(D) + CSZA & Use Gamma(D) instead of Gamma(A); accompanying controller variants. 69-reader/nine-fold or ten-fold allocations. \\
Known-zero source returns & PA194A; PA194F; PA194J; PA194ZR; PA194R; PA194CSZ = PA194X; PA194CZ; PA194HA; PA194HAZ; CSZA final-register 3; CSZA final-register 5; CSZA final-register 6 & Replace selected jumps or returns by decrements of registers proved zero. Changes recurring routing and, in R, decoder association. \\
Physical register placement & PA194-AD; PA194F\_AD; PA194R\_AD; PA194TDY; PA194P5; PA194HX & Preserve logical commands while exchanging register positions and routers. Two-, three- and five-register arrangements. \\
Decoder topology & PA194B; PA1944L & Preserve the original source program. Dedicated history selector; four-increment reassociation plus helper/fold. \\
AX11 interface & PA194AX11 & Push the guard and fall through the full MP operand reader: four operands, eleven entries. Reuses existing Push and reader machinery. \\
AX11 with cache-only guard & PA194AX11H; PA194AX11QH & Four-operand AX11 with H guard, with original or Q distribution. One source increment removed; same 194-state budget. \\
AX11 with changed coding & PA194AX11Q6; PA194AX11QH5 & AX11 and Q with coefficient six/sum guard or five/cache guard. 135 fixed plus 67 readers minus eight folds. \\
AX11 with symmetric sum guard & PA194AX11S & Exactly the AX11 sum-guard command semantics. Symmetric source tail; three fewer source instructions. \\
Explicit generalization binder & PA194G; PA194AX11G & The first AX operand c controls U(c,h), instead of inherited offset d; ten- or eleven-entry AX. Q emitter and no-push MP; original sum guard. \\
Explicit binder with H guard & PA194GH; PA194GHS & Explicit binder and ten-entry AX with cache-only Shift. Ordinary versus symmetric comparator. \\
AX11QH decoder alternatives & PA194AX11QH4L; PA194AX11QH4R & Exactly the existing four-operand AX11QH source. Two four-increment reassociations with different relative time bounds. \\
\bottomrule
\end{longtable}
